% New post-checkpoint05 contribution.  ID A.4.23 reserved by root.
\subsection{Pure orthogonal corrections preserving the entire lower comparison}

\begin{result}{A.4.23}{Proved}{Every perfect orthogonal correction has an actual pure lift with trivial lower image}
\label{prop:pure-orthogonal-corrections}
Let \(\Pi=\Pi_{1,9}\), with its nine named peripheral conjugacy
classes and its specified quaternion quotient from A.4.15.
Let \(\mathcal G\) be the group of actual based,
orientation-preserving geometric automorphisms of \(\Pi\)
which fix all nine labels and induce the identity automorphism
of this \(Q_8\) quotient. Its action on the specified multiplicity
space is the homomorphism
\[
 \rho:\mathcal G\longrightarrow\mathrm{SO}_{18}^-(3).
 \tag{PC.1}
\]
This homomorphism is surjective.

Retain the actual joint finite lower map
\[
 \begin{split}
 \beta:\mathcal G&\longrightarrow\mathcal B,\\
 g&\longmapsto\bigl([g|_{\Pi/(\gamma_3\Pi\Pi^3)}],\,
                    g|_{H_1(\Pi,\mathbb Z/9)},\\
   &\hspace{42mm}g|_{H_1(\overline\Pi,\mathbb Z/12)}\bigr),
 \end{split}
 \tag{PC.2}
\]
where the first coordinate is an outer automorphism, the second
is the literal homology action in the marked rank-ten basis,
the third is the signed ordinary compact-period matrix, and
\(\mathcal B\) denotes the actual image of this joint map.
Here \(\overline\Pi\) is the quotient of \(\Pi\) by the normal
closure of its nine peripheral elements, so the third coordinate
is induced by this specified compactification map.
Then
\[
                  \rho(\ker\beta)\supseteq\Omega_{18}^-(3).
 \tag{PC.3}
\]
More precisely, for every \(\omega\in\Omega_{18}^-(3)\) there
is an actual finite word in pure quaternion-kernel mapping
classes whose multiplicity matrix is \(\omega\) and all three
coordinates of (PC.2) are the identity. The correcting word
also fixes all puncture labels and the full specified quaternion
quotient, by its membership in \(\mathcal G\).
\end{result}
\begin{proof}
We first verify surjectivity on the stated pure domain.
In the fifteen actual generators of A.4.15, the root matrix
with zero-based index forty-nine among the quaternion-kernel
Schreier generators has the literal source word
\[
               (3,5,-3)
       =\mathrm{push}_1\,\sigma_3\,\mathrm{push}_1^{-1}.
 \tag{PC.4}
\]
The point push fixes all labels, and the displayed disk half
twist interchanges two of them. Thus (PC.4) has a transposition
as its puncture permutation. It lies in the quaternion kernel
already, and its square is pure. Its multiplicity image is
the certified root element \(T_e(v)\), so the image of its
square is exactly \(T_e(2v)\), by the additive root formula.

Each of the sixteen certified conjugates of this word has
a transposition as its label permutation. Their squares
therefore lie in \(\mathcal G\), and their parameters are
twice the sixteen certified basis parameters of
\(e^\perp/\langle e\rangle\). They still span that space.
Consequently \(\rho(\mathcal G)\) contains the entire first
radical \(R_e\). The pure subgroup is normal in the full
geometric group, and the quaternion kernel is normal as well.
Conjugating these pure words by the actual opposite-line
generator, with zero-based index eleven, therefore gives
pure words in the same quaternion kernel. Their image is
the entire opposite radical. A.4.16 now gives
\(\Omega_{18}^-(3)\subseteq\rho(\mathcal G)\).

There is also a pure word in the other spinor coset.
The quaternion-kernel Schreier generator with zero-based
index fifteen is the literal word \((1,1,-13)\).
Its puncture permutation is a nontrivial translation of
order three; its certified spinor is the nonsquare class.
Its cube is therefore pure and has the same nonsquare
spinor, because \(2^3=2\) in \(\mathbb F_3^\times\).
It belongs to \(\mathcal G\). Since the target of (PC.1)
is already contained in \(\mathrm{SO}\), and its kernel
under spinor is \(\Omega\), this proves surjectivity of
(PC.1).

All these source assertions are independently checked in
\texttt{check\_main\_pure\_radicals.py}. It computes the
nine-label permutations from the original labelled covering,
expands every conjugating word to the fifteen source generators,
checks that each of the thirty-two squared/conjugated words
is pure, and verifies their actual root matrices. The expanded
words and all label permutations are retained. The additional
cube is checked against its recorded ordinary-reflection
factorization.

We record explicitly the perfectness used in the lower-image
argument. Theorem11.47 of \cite[pp.~161--162]{GhostTaylorClassical}
asserts that \(\Omega_{18}^-(3)\) is perfect; all its exceptions
have dimension three or four. In the present coordinates the
characteristic-three part of its proof is short. For a
nonzero singular \(e\), choose an opposite \(f\), and represent
a root parameter \(v\) in the nondegenerate space
\(W=\langle e,f\rangle^\perp\), of dimension sixteen.
For \(q(v)\ne0\), the nondegenerate space
\(W\cap v^\perp\), of dimension fifteen, contains an \(x\)
with \(q(x)=q(v)\). Indeed diagonalize two of its directions:
one has this norm, or both have the other nonzero norm, in
which case their sum has the required norm over \(\mathbb F_3\).
Let \(h\) be minus the identity on the orthogonal plane
\(\langle v,x\rangle\), and the identity on its orthogonal
complement. It is the product of the two reflections in
\(v,x\); its determinant is one and its spinor is
\(q(v)q(x)=1\). Thus \(h\in\Omega\), and
\[
                  [h,T_e(v)]=T_e(-2v)=T_e(v).
 \tag{PC.5}
\]
Both entries of this commutator lie in \(\Omega\).
For a singular nonzero parameter \(v\in W\), the space
\(W\cap v^\perp\) has radical \(\langle v\rangle\)
and a nondegenerate quotient of dimension fourteen, so
it contains an anisotropic \(w\). Both \(w\) and \(v-w\)
are anisotropic, since their norms are equal and nonzero.
The additive root formula expresses \(T_e(v)\) as the
product of their two root elements. Thus all root elements
lie in \([\Omega,\Omega]\), and they generate \(\Omega\).
This proves perfectness in the exact dimension and field
used here. It also gives
\[
             \mathrm{SO}_{18}^-(3)^{(d)}=\Omega_{18}^-(3)
                          \qquad(d\ge1),
 \tag{PC.6}
\]
because \(\mathrm{SO}/\Omega\) has order two and \(\Omega\)
is perfect.

The image \(\mathcal B\) of (PC.2) is solvable, with
derived length at most five. For the first coordinate,
A.4.9 gives an elementary abelian kernel of dimension
ninety-one over the homology image, and A.4.8 identifies
that homology image as an extension of an elementary
abelian group of dimension sixteen by
\(\mathrm{SL}_2(\mathbb F_3)\). The latter has the
derived series
\(\mathrm{SL}_2(\mathbb F_3),Q_8,C_2,1\).
These two abelian extensions give derived length at
most five. The same bound holds for the second coordinate:
the rank-ten pure homology matrices modulo nine have
an abelian lower-left block over
\(\mathrm{SL}_2(\mathbb Z/9)\), whose kernel modulo
three is the additive trace-zero matrix space, an abelian
group. For the third coordinate, reduction modulo four
has abelian kernel over \(\mathrm{SL}_2(\mathbb F_2)\),
and reduction modulo three is
\(\mathrm{SL}_2(\mathbb F_3)\). Thus this signed-period
group is solvable of derived length at most three.
A subgroup of the product has derived length at most
five, proving the assertion for the actual joint image
\(\mathcal B\).

Now form the actual finite joint image
\[
 D=(\rho,\beta)(\mathcal G)
     \le\ \mathrm{SO}_{18}^-(3)\times\mathcal B.
 \tag{PC.7}
\]
Projection \(\pi:D\to\mathrm{SO}_{18}^-(3)\) is
surjective by the pure word construction above.
For a surjective group homomorphism, images of derived
subgroups satisfy
\[
                  \pi(D^{(d)})=\pi(D)^{(d)}.
 \tag{PC.8}
\]
For \(d=1\), every commutator of two image elements
is the image of the commutator of chosen preimages,
which proves both containments. Induction proves the
display for every \(d\).
Projection of \(D^{(5)}\) to \(\mathcal B\) is trivial,
whereas (PC.6)--(PC.8) give
\[
            \pi(D^{(5)})=\Omega_{18}^-(3).
 \tag{PC.9}
\]
Every element of \(D^{(5)}\) is the image of a product
of iterated commutators of actual elements of
\(\mathcal G\). These are finite geometric words,
and their three lower coordinates are simultaneously
the identity. This proves (PC.3) with precisely the
source and lower maps stated, and proves the stronger
word assertion.
\end{proof}
